% Generated draft: requires mathematical and source review.
\documentclass[11pt]{article}
\usepackage[margin=25mm]{geometry}
\usepackage{fontspec}
\setmainfont{Latin Modern Roman}
\usepackage{amsmath,amssymb,mathtools,unicode-math}
\setmathfont{Latin Modern Math}
\usepackage{enumitem,needspace,xurl,hyperref}
\setlength{\parindent}{0pt}
\setlength{\parskip}{5pt}
\setlength{\emergencystretch}{3em}
\newcommand{\N}{\mathbb{N}}
\newcommand{\Z}{\mathbb{Z}}
\newcommand{\Q}{\mathbb{Q}}
\newcommand{\R}{\mathbb{R}}
\newcommand{\C}{\mathbb{C}}
\newcommand{\F}{\mathbb{F}}
\newcommand{\A}{\mathbb{A}}
\newcommand{\PP}{\mathbb{P}}
\newcommand{\E}{\mathbb{E}}
\newcommand{\eps}{\varepsilon}
\newcommand{\dd}{\,\mathrm{d}}
\begin{document}
\section*{UnsolvedMath ID: 30006628}
\textbf{Source problem code:} catalogue entry
\subsection{Definitions and mathematical statement}
Let \(G=\mathop{\ast}_{i\in I}G_i\) be a free product of nontrivial groups, with \(|I|\geq 3\). The supplied record does not specify whether infinitely many factors are allowed or what additional hypotheses the intended random model imposes on the factors.

A group \(H\) has Kazhdan's property \((T)\) if every unitary representation \(\pi:H\to\mathcal U(\mathcal H)\) on a complex Hilbert space that has almost invariant unit vectors has a nonzero invariant vector. Here almost invariant unit vectors means that for every finite subset \(F\subseteq H\) and every \(\varepsilon>0\), there exists \(v\in\mathcal H\) with \(\|v\|=1\) and \(\|\pi(h)v-v\|<\varepsilon\) for all \(h\in F\). An invariant vector satisfies \(\pi(h)v=v\) for every \(h\in H\).

For a random set of relators \(R_{L,d}\subseteq G\), write
\[
Q_{L,d}=G/\langle\!\langle R_{L,d}\rangle\!\rangle_G,
\]
where the denominator is the normal subgroup of \(G\) generated by \(R_{L,d}\). In this notation, \(L\) denotes the intended model's asymptotic size parameter and \(d\) its density parameter; these are placeholders because the supplied record does not define the sampling law or the asymptotic regime.

\textbf{Question, conditional on specifying the intended free-product density model.} For every fixed free product with at least three nontrivial factors admitted by that model, and every fixed admissible density \(d>1/3\), is it true that
\[
\lim_{L\to\infty}\mathbb P_{L,d}\bigl(Q_{L,d}\text{ has Kazhdan's property }(T)\bigr)=1?
\]
Here \(\mathbb P_{L,d}\) must be the probability law prescribed by the intended model, and the limit must follow that model's permitted parameter sequence. The inequality is strict; the question makes no assertion at \(d=1/3\).

The supplied record does not specify the relator alphabet, word conventions, number or distribution of relators, allowable density range, or limiting sequence. Consequently, the display records the intended quantifiers and target property but is not a fully specified probabilistic statement until those model details are recovered. The original excerpt asks whether ``this'' remains true; the supplied cleaned statement identifies the target as property \((T)\), but the antecedent passage is absent.
\subsection{Short English statement}
The question asks whether random quotients of every admissible free product of at least three nontrivial groups have Kazhdan's property \((T)\) with probability tending to one at each fixed density \(d>1/3\). The precise random model and its hypotheses require clarification. The supplied registry labels the problem open; this draft does not independently verify that status.
\subsection{Sources}
\begin{itemize}
\item Local UnsolvedMath registry, supplied complete source record: permanent ID \(30006628\), exact source problem code \texttt{catalogue entry}, title \emph{Property T for Random Free-Product Quotients above One-Third Density}. The record gives proposed year 2026 but names no proposer. Its statement status \texttt{corrected\_verified} and accompanying claim of typographical repairs are supplied registry metadata, not independent verification here.
\item Supplied report attribution: \emph{Median Geometry and Applications} (2026), Oberwolfach Reports, pp.~487--540, DOI \texttt{10.4171/owr/2026/8}. \url{https://doi.org/10.4171/owr/2026/8}. The background describes Oberwolfach Reports as workshop reports containing open-problem discussions; no precise page for this question is supplied.
\item Supplied literature source, labeled ``OWR report'': \url{https://ems.press/content/serial-article-files/53603}.
\item Supplied literature source, labeled ``Einstein--Krishna--Montee--Ng--Steenbock preprint'': \url{https://arxiv.org/abs/2502.08630}. The supplied background describes this as a 2025 random-quotients paper.
\item Literature assessment dated exactly \texttt{2026-08-22}: the supplied curation labels the problem open, says the question for at least three factors is explicitly left open in the 2026 report and the 2025 random-quotients paper, and reports that no resolution was found. This is supplied dated curation derived from the OWR record, not an independently verified literature assessment. The source itself requests independent verification of time-sensitive status and cited claims.
\end{itemize}
\end{document}
